% GENERATED FILE. Edit canonical lessons, then run npm run generate:books.
% canonical-source-hash: fnv1a64:c425f779ed9d8229
% canonical-lessons: HI-C459-sneh, HI-C459-mamta, HI-C459-lar, HI-C459-cumna, HI-C459-puckarna, HI-R459-review-452-459

\chapter{Affection}
\label{ch:hi-affection}
\hlchaptermodality{459}

\begin{chapteropening}
\textbf{By the end of this chapter:} \emph{I can say how someone shows their love.}

Complete the last lesson of chapter 459: I can say how someone shows their love.
\end{chapteropening}

\section[\texorpdfstring{\hi{स्नेह} (sneh)}{स्नेह (sneh)}]{\hi{स्नेह} (sneh) — affection; fondness}

\label{lesson:HI-C459-sneh}



\begin{quote}
Before the new one: say the Hindi for to writhe, then the Hindi for an assurance.
\end{quote}

\subsection*{You'll want to know: \hi{स्नेह}}
\textbf{\hi{स्नेह}} — \emph{sneh} — \textquotedblleft{}affection; fondness\textquotedblright{}.

\begin{quote}
\hi{दादी} \hi{सब} \hi{बच्चों} \hi{से} \hi{बहुत} \hi{स्नेह} \hi{करती} \hi{हैं।}
\end{quote}

\begin{quote}
\emph{dādī sab baccoṁ se bahut sneh kartī haiṁ.}
\end{quote}

\begin{quote}
Grandmother is very fond of all the children.
\end{quote}

Masculine.

\begin{cousinweb}
\textbf{\hi{स्नेह}} (\emph{sneh}) is a tatsama, borrowed from Sanskrit \emph{sneha}, first \textquotedblleft{}oiliness\textquotedblright{} and then \textquotedblleft{}tenderness, love, fondness\textquotedblright{} (Monier-Williams).
\end{cousinweb}

\begin{grammarlens}[title={What you've built}]
One more word for describing what you lived and felt.
\end{grammarlens}

\subsection*{Your turn}
\begin{itemize}
\raggedright
  \item \emph{Say it:} \emph{sneh}
  \item \emph{Say it:} \emph{sneh}, once more
  \item \emph{Say it:} \emph{dādī sab baccoṁ se bahut sneh kartī haiṁ.} — Grandmother is very fond of all the children.
  \item \emph{Recall:} say the Hindi for to writhe, then the Hindi for an assurance, then say \emph{sneh} again
\end{itemize}

\begin{tcolorbox}[breakable,colback=teal!4,colframe=teal!35!black,title={Before you move on}]
What does \hi{स्नेह} mean? (\textquotedblleft{}Affection; fondness\textquotedblright{}.) Say it once more.
\end{tcolorbox}
\section[\texorpdfstring{\hi{ममता} (mamtā)}{ममता (mamtā)}]{\hi{ममता} (mamtā) — a mother's love}

\label{lesson:HI-C459-mamta}



\begin{quote}
Before the new one: say the Hindi for an assurance, then the Hindi for affection.
\end{quote}

\subsection*{You'll want to know: \hi{ममता}}
\textbf{\hi{ममता}} — \emph{mamtā} — \textquotedblleft{}a mother's love\textquotedblright{}.

\begin{quote}
\hi{माँ} \hi{की} \hi{ममता} \hi{की} \hi{कोई} \hi{क़ीमत} \hi{नहीं} \hi{होती।}
\end{quote}

\begin{quote}
\emph{māṁ kī mamtā kī koī qīmat nahīṃ hotī.}
\end{quote}

\begin{quote}
No price can be put on a mother's love.
\end{quote}

Feminine.

\begin{cousinweb}
\textbf{\hi{ममता}} (\emph{mamtā}) is a tatsama, borrowed from Sanskrit \emph{mamatā}, \textquotedblleft{}the state of 'mine'\textquotedblright{}, built on \emph{mama}, \textquotedblleft{}mine\textquotedblright{} (Monier-Williams). Hindi turned the sense of ownership into the tender love of a mother.
\end{cousinweb}

\begin{grammarlens}[title={What you've built}]
One more word for describing what you lived and felt.
\end{grammarlens}

\subsection*{Your turn}
\begin{itemize}
\raggedright
  \item \emph{Say it:} \emph{mamtā}
  \item \emph{Say it:} \emph{mamtā}, once more
  \item \emph{Say it:} \emph{māṁ kī mamtā kī koī qīmat nahīṃ hotī.} — No price can be put on a mother's love.
  \item \emph{Recall:} say the Hindi for an assurance, then the Hindi for affection, then say \emph{mamtā} again
\end{itemize}

\begin{tcolorbox}[breakable,colback=teal!4,colframe=teal!35!black,title={Before you move on}]
What does \hi{ममता} mean? (\textquotedblleft{}A mother's love\textquotedblright{}.) Say it once more.
\end{tcolorbox}
\section[\texorpdfstring{\hi{लाड़} (lāṛ)}{लाड़ (lāṛ)}]{\hi{लाड़} (lāṛ) — fond pampering; fondness}

\label{lesson:HI-C459-lar}



\begin{quote}
Before the new one: say the Hindi for affection, then the Hindi for a mother's love.
\end{quote}

\subsection*{You'll want to know: \hi{लाड़}}
\textbf{\hi{लाड़}} — \emph{lāṛ} — \textquotedblleft{}fond pampering; fondness\textquotedblright{}.

\begin{quote}
\hi{नाना} \hi{के} \hi{लाड़} \hi{में} \hi{मीरा} \hi{का} \hi{बचपन} \hi{बीता।}
\end{quote}

\begin{quote}
\emph{nānā ke lāṛ meṁ Mīrā kā bacpan bītā.}
\end{quote}

\begin{quote}
Mira's childhood was spent in her grandfather's fond care.
\end{quote}

Masculine.

\begin{cousinweb}
\textbf{\hi{लाड़}} (\emph{lāṛ}) is a tadbhava: CDIAL lists Hindi \emph{lāṛ}, \textquotedblleft{}affection\textquotedblright{}, under \emph{*lāḍya}, a form it reconstructs (CDIAL 11013).
\end{cousinweb}

\begin{grammarlens}[title={What you've built}]
One more word for describing what you lived and felt.
\end{grammarlens}

\subsection*{Your turn}
\begin{itemize}
\raggedright
  \item \emph{Say it:} \emph{lāṛ}
  \item \emph{Say it:} \emph{lāṛ}, once more
  \item \emph{Say it:} \emph{nānā ke lāṛ meṁ Mīrā kā bacpan bītā.} — Mira's childhood was spent in her grandfather's fond care.
  \item \emph{Recall:} say the Hindi for affection, then the Hindi for a mother's love, then say \emph{lāṛ} again
\end{itemize}

\begin{tcolorbox}[breakable,colback=teal!4,colframe=teal!35!black,title={Before you move on}]
What does \hi{लाड़} mean? (\textquotedblleft{}Fond pampering; fondness\textquotedblright{}.) Say it once more.
\end{tcolorbox}
\section[\texorpdfstring{\hi{चूमना} (cūmnā)}{चूमना (cūmnā)}]{\hi{चूमना} (cūmnā) — to kiss}

\label{lesson:HI-C459-cumna}



\begin{quote}
Before the new one: say the Hindi for a mother's love, then the Hindi for fond pampering.
\end{quote}

\subsection*{You'll want to know: \hi{चूमना}}
\textbf{\hi{चूमना}} — \emph{cūmnā} — \textquotedblleft{}to kiss\textquotedblright{}.

\begin{quote}
\hi{माँ} \hi{ने} \hi{बच्चे} \hi{का} \hi{माथा} \hi{चूमा।}
\end{quote}

\begin{quote}
\emph{māṁ ne bacce kā māthā cūmā.}
\end{quote}

\begin{quote}
The mother kissed the child's forehead.
\end{quote}

Past, with \hi{ने} (he or she …): \hi{उसने} \hi{माथा} \hi{चूमा}, \emph{usne māthā cūmā}.

\begin{cousinweb}
\textbf{\hi{चूमना}} (\emph{cūmnā}) is a tadbhava: CDIAL lists Hindi \emph{cūmnā}, \textquotedblleft{}to kiss\textquotedblright{}, under Sanskrit \emph{cúmbati}, \textquotedblleft{}kisses\textquotedblright{} (CDIAL 4870).
\end{cousinweb}

\begin{grammarlens}[title={What you've built}]
One more word for describing what you lived and felt.
\end{grammarlens}

\subsection*{Your turn}
\begin{itemize}
\raggedright
  \item \emph{Say it:} \emph{cūmnā}
  \item \emph{Say it:} \emph{cūmnā}, once more
  \item \emph{Say it:} \emph{māṁ ne bacce kā māthā cūmā.} — The mother kissed the child's forehead.
  \item \emph{Recall:} say the Hindi for a mother's love, then the Hindi for fond pampering, then say \emph{cūmnā} again
\end{itemize}

\begin{tcolorbox}[breakable,colback=teal!4,colframe=teal!35!black,title={Before you move on}]
What does \hi{चूमना} mean? (\textquotedblleft{}To kiss\textquotedblright{}.) Say it once more.
\end{tcolorbox}
\section[\texorpdfstring{\hi{पुचकारना} (puckārnā)}{पुचकारना (puckārnā)}]{\hi{पुचकारना} (puckārnā) — to soothe with soft words}

\label{lesson:HI-C459-puckarna}



\begin{quote}
Before the new one: say the Hindi for fond pampering, then the Hindi for to kiss.
\end{quote}

\subsection*{You'll want to know: \hi{पुचकारना}}
\textbf{\hi{पुचकारना}} — \emph{puckārnā} — \textquotedblleft{}to soothe with soft words\textquotedblright{}.

\begin{quote}
\hi{बच्चा} \hi{रो} \hi{रहा} \hi{था}, \hi{तो} \hi{दादी} \hi{ने} \hi{उसे} \hi{पुचकारा।}
\end{quote}

\begin{quote}
\emph{baccā ro rahā thā, to dādī ne use puckārā.}
\end{quote}

\begin{quote}
The child was crying, so my grandmother soothed him with soft words.
\end{quote}

Past, with \hi{ने} (he or she …): \hi{उसने} \hi{बच्चे} \hi{को} \hi{पुचकारा}, \emph{usne bacce ko puckārā}.

\begin{cousinweb}
\textbf{\hi{पुचकारना}} (\emph{puckārnā}) is a tadbhava: CDIAL lists Hindi \emph{puckārnā}, \textquotedblleft{}to caress\textquotedblright{}, under \emph{*pucca}, a form it reconstructs (CDIAL 8248). It is the soft clicking and cooing used to calm a child or an animal.
\end{cousinweb}

\begin{grammarlens}[title={What you've built}]
One more word for describing what you lived and felt.
\end{grammarlens}

\subsection*{Your turn}
\begin{itemize}
\raggedright
  \item \emph{Say it:} \emph{puckārnā}
  \item \emph{Say it:} \emph{puckārnā}, once more
  \item \emph{Say it:} \emph{baccā ro rahā thā, to dādī ne use puckārā.} — The child was crying, so my grandmother soothed him with soft words.
  \item \emph{Recall:} say the Hindi for fond pampering, then the Hindi for to kiss, then say \emph{puckārnā} again
\end{itemize}

\begin{tcolorbox}[breakable,colback=teal!4,colframe=teal!35!black,title={Before you move on}]
What does \hi{पुचकारना} mean? (\textquotedblleft{}To soothe with soft words\textquotedblright{}.) Say it once more.
\end{tcolorbox}
\section[(dialogue)]{Looking back: from Naming a Feeling to Affection}

\label{lesson:HI-R459-review-452-459}



\begin{quote}
Say the Hindi for a feeling and an awareness.
\end{quote}

\subsection*{Your turn}
Say these aloud:

\begin{itemize}
\raggedright
  \item a feeling, an awareness, a deep feeling, to feel, to well up
  \item joy, zeal, a smile, elation, pleased
  \item taken aback, surprise, astonishment, strange, a marvel
  \item grief, disappointed, dejected, to fret, to hurt someone's feelings
  \item cross, wrath, to get irritated, to fly into a rage, to get fed up
  \item dread, to panic, to frighten, cowardly, to shudder
  \item stress, concern, a burden, to writhe, an assurance
  \item affection, a mother's love, fond pampering, to kiss, to soothe with soft words
\end{itemize}

\begin{tcolorbox}[breakable,colback=teal!4,colframe=teal!35!black,title={Before you move on}]
A feeling? (\textbf{\hi{भावना}}, \emph{bhāvnā}.) An awareness? (\textbf{\hi{एहसास}}, \emph{ehsās}.) A deep feeling? (\textbf{\hi{अनुभूति}}, \emph{anubhūti}.) To feel? (\textbf{\hi{महसूस} \hi{करना}}, \emph{mahsūs karnā}.) To well up? (\textbf{\hi{उमड़ना}}, \emph{umaṛnā}.) Joy? (\textbf{\hi{आनंद}}, \emph{ānand}.) Zeal? (\textbf{\hi{उत्साह}}, \emph{utsāh}.) A smile? (\textbf{\hi{मुस्कान}}, \emph{muskān}.) Elation? (\textbf{\hi{उल्लास}}, \emph{ullās}.) Pleased? (\textbf{\hi{प्रसन्न}}, \emph{prasann}.) Taken aback? (\textbf{\hi{चकित}}, \emph{cakit}.) Surprise? (\textbf{\hi{आश्चर्य}}, \emph{āścarya}.) Astonishment? (\textbf{\hi{अचंभा}}, \emph{acambhā}.) Strange? (\textbf{\hi{अजीब}}, \emph{ajīb}.) A marvel? (\textbf{\hi{कमाल}}, \emph{kamāl}.) Grief? (\textbf{\hi{ग़म}}, \emph{ġam}.) Disappointed? (\textbf{\hi{निराश}}, \emph{nirāś}.) Dejected? (\textbf{\hi{मायूस}}, \emph{māyūs}.) To fret? (\textbf{\hi{कुढ़ना}}, \emph{kuṛhnā}.) To hurt someone's feelings? (\textbf{\hi{दुखाना}}, \emph{dukhānā}.) Cross? (\textbf{\hi{नाराज़}}, \emph{nārāz}.) Wrath? (\textbf{\hi{क्रोध}}, \emph{krodh}.) To get irritated? (\textbf{\hi{चिढ़ना}}, \emph{ciṛhnā}.) To fly into a rage? (\textbf{\hi{बिफरना}}, \emph{bipharnā}.) To get fed up? (\textbf{\hi{खीझना}}, \emph{khījhnā}.) Dread? (\textbf{\hi{भय}}, \emph{bhay}.) To panic? (\textbf{\hi{घबराना}}, \emph{ghabrānā}.) To frighten? (\textbf{\hi{डराना}}, \emph{ḍarānā}.) Cowardly? (\textbf{\hi{कायर}}, \emph{kāyar}.) To shudder? (\textbf{\hi{सिहरना}}, \emph{siharnā}.) Stress? (\textbf{\hi{तनाव}}, \emph{tanāv}.) Concern? (\textbf{\hi{फ़िक्र}}, \emph{fikr}.) A burden? (\textbf{\hi{बोझ}}, \emph{bojh}.) To writhe? (\textbf{\hi{तड़पना}}, \emph{taṛapnā}.) An assurance? (\textbf{\hi{आश्वासन}}, \emph{āśvāsan}.) Affection? (\textbf{\hi{स्नेह}}, \emph{sneh}.) A mother's love? (\textbf{\hi{ममता}}, \emph{mamtā}.) Fond pampering? (\textbf{\hi{लाड़}}, \emph{lāṛ}.) To kiss? (\textbf{\hi{चूमना}}, \emph{cūmnā}.) To soothe with soft words? (\textbf{\hi{पुचकारना}}, \emph{puckārnā}.)
\end{tcolorbox}
